\documentclass[12pt]{extarticle}
\usepackage[utf8]{inputenc}
\usepackage{amsmath}
\usepackage{amssymb}
\usepackage{graphicx}
\usepackage{geometry}
\usepackage{hyperref}
\usepackage[T1]{fontenc}
\usepackage[charter]{mathdesign}
\usepackage{tcolorbox}
\usepackage{array}
\usepackage{enumitem}
\usepackage[table]{xcolor}
\usepackage{enumitem}
\usepackage{tikz}
\usepackage{fancyhdr}
\usepackage{titlesec}
\usepackage{microtype}
\usepackage{setspace}
\usepackage{tabularx}
\usepackage{booktabs}
\usepackage{lastpage}
\usepackage{listings}
\usepackage{etoolbox}
\usetikzlibrary{calc}
\usetikzlibrary{decorations.markings}

\geometry{
 a4paper,
 left=18mm,
 right=18mm,
 top=20mm,
 bottom=22mm,
 headheight=10mm,
}

\definecolor{headerNavy}{HTML}{1B3A5C}
\definecolor{header}{HTML}{b1caf5}
\definecolor{Narangi}{HTML}{EF6820}
\definecolor{accentGold}{HTML}{C5961D}
\definecolor{accentBlue}{HTML}{2E6DA4}
\definecolor{softGray}{HTML}{6B7280}
\definecolor{darkText}{HTML}{1A202C}
\definecolor{lightBg}{HTML}{F7F8FA}
\definecolor{borderMed}{HTML}{CBD5E1}
\definecolor{borderLight}{HTML}{E2E8F0}
\definecolor{pastelBlue}{HTML}{DBEAFE}
\definecolor{pastelMint}{HTML}{D1FAE5}
\definecolor{pastelLavender}{HTML}{E9D5FF}
\definecolor{cdOrange}{HTML}{E8650A}
\definecolor{ruleColor}{HTML}{94A3B8}
\definecolor{brandAccent}{HTML}{06B6D4}
\definecolor{mycolor}{HTML}{1669b3}

\definecolor{solGreen}{HTML}{16A34A}
\definecolor{solBg}{HTML}{F0FDF4}

% Auto-green "Correct Answer:" + auto-boxed "Solution:" via preamble-only \textbf redefinition
\newif\ifinsolutionbox
\let\oldtextbf\textbf
\renewcommand{\textbf}[1]{%
  \def\temp{#1}%
  \def\tempcomp{Correct Answer:}%
  \def\tempsol{Solution:}%
  \ifx\temp\tempcomp
    {\color{solGreen}\oldtextbf{#1}}%
  \else\ifx\temp\tempsol
    \begin{tcolorbox}[enhanced, breakable, colback=solBg, colframe=solGreen!60, boxrule=0.5pt, leftrule=3pt, arc=2pt, left=10pt, right=10pt, top=6pt, bottom=6pt]%
    \global\insolutionboxtrue
    {\color{solGreen}\oldtextbf{#1}}%
  \else
    \oldtextbf{#1}%
  \fi\fi
}

% Auto-close the solution box right before each quicktipbox
\BeforeBeginEnvironment{quicktipbox}{%
  \ifinsolutionbox
    \end{tcolorbox}%
    \global\insolutionboxfalse
  \fi
}

\setstretch{1.45}
\setlength{\parindent}{0pt}
\setlength{\parskip}{3pt}

\pagestyle{fancy}
\fancyhf{}
\renewcommand{\headrulewidth}{0pt}
\renewcommand{\footrulewidth}{0.4pt}
\renewcommand{\footrule}{\hbox to\headwidth{\color{borderMed}\leaders\hrule height
\footrulewidth\hfill}}

\fancyfoot[L]{\scriptsize\color{softGray} AIIMS Paramedical 2026}
\fancyfoot[C]{\small\color{softGray} Page \textbf{\thepage}\ of \textbf{\pageref{LastPage}}}
\fancyfoot[R]{\scriptsize\color{softGray} \copyright\ Collegedunia.com}

\tcbuselibrary{skins, breakable}

% Environment for quick tip boxes (named quicktipbox for Python script compatibility)
\newtcolorbox{quicktipbox}{
  enhanced,
  colback=brandAccent!5,
  colframe=brandAccent,
  boxrule=0.5pt,
  leftrule=3pt,
  arc=2pt,
  left=10pt, right=10pt, top=6pt, bottom=6pt,
  fontupper=\small\color{darkText},
  before upper={\textbf{\color{brandAccent}Quick Tip:}\hspace{6pt}}
}

\newtcolorbox{sectionheader}[1][]{
  enhanced,
  colback=headerNavy!8,
  colframe=headerNavy!40,
  boxrule=0.6pt,
  arc=4pt,
  left=10pt, right=10pt, top=8pt, bottom=8pt,
  fontupper=\large\bfseries\color{headerNavy},
  halign=center,
  #1
}

\newtcolorbox{solutionbox}{
  breakable,
  enhanced,
  colback=solBg,
  colframe=solGreen,
  boxrule=0.8pt,
  arc=4pt,
  left=10pt, right=10pt, top=8pt, bottom=8pt
}

\newcommand{\qnum}[1]{
  \textbf{\color{headerNavy}#1.}\hspace{3pt}
}

\newcommand{\questiondivider}{
  \vspace{5pt}
  {\centering\color{borderLight}\rule{0.5\textwidth}{0.3pt}\par}
  \vspace{5pt}
}

\begin{document}
\thispagestyle{empty}
\noindent
\begin{tikzpicture}[remember picture, overlay]
  % 1. Main Background
  \fill[pastelBlue!25]
    ([yshift=0mm]current page.north west) rectangle 
    ([yshift=-48mm]current page.north east);

  % 3. Bottom Accent Strip
  \fill[Narangi!60]
    ([yshift=-48mm]current page.north west) rectangle
    ([yshift=-49.5mm]current page.north east);

  % 4. Logo Halo (White circle behind logo to make it stand out)
  \fill[white, opacity=0.8] 
    ([xshift=-29mm, yshift=-24mm]current page.north east) circle (19mm);

  % 5. Right: Logo
  \node[anchor=north east]
    at ([xshift=-12mm, yshift=-9mm]current page.north east)
    {\includegraphics[width=0.19\linewidth]{AIIMS.png}};

  % 6. Text Content
  % Line 1: Main Title
  \node[anchor=north west, text=headerNavy, font=\fontsize{22}{28}\selectfont\bfseries] (Title)
    at ([xshift=18mm, yshift=-12mm]current page.north west)
    {AIIMS Paramedical 2026 };

  % Small divider line

  % Line 2: PCM Subtitle
  \node[anchor=north west, text=black!85, font=\fontsize{17}{19}\selectfont\bfseries] (Subtitle)
    at ([xshift=18mm, yshift=-25mm]current page.north west)
    {Question Paper (Memory-Based) with Solutions};

  % Line 3: Authority
  \node[anchor=north west, text=black!60, font=\footnotesize] (Authority)
    at ([xshift=18mm, yshift=-33mm]current page.north west)
    {Conducted by AIIMS Delhi };

\end{tikzpicture}

\vspace{25mm}

\begin{tcolorbox}[
    colframe=Narangi!60,
    colback=white,
    arc=3mm,
    boxrule=1.5pt,
    left=15pt, right=15pt,
    top=10pt, bottom=10pt
]
    {\large\bfseries\color{headerNavy} General Instructions}\\[2pt]
    {\color{Narangi!60}\rule{4.5cm}{1.5pt}}

\begin{enumerate}[leftmargin=20pt, itemsep=4pt, label={\textbf{\color{headerNavy}(\roman*)}}]
    \item Duration of the exam is 90 minutes.
    \item This question paper consists of 90 questions. 
    \item Question Paper is divided into 3 sections- Physics, Chemistry and Mathematics or Biology having 30 questions each.
    \item Each question carries 1 mark and their is negative marking of 1/3 for every incorrect answer.
\end{enumerate}
\end{tcolorbox}

\vspace{6mm}

\begin{sectionheader}[colback=pastelBlue!25, colframe=pastelBlue!55]
Biology
\end{sectionheader}
\vspace{4mm}


\noindent \textbf{1.}  
% Question text  
\textbf{Oxytocin and Vasopressin are secreted by:} \\  

% Option  
(A) Anterior pituitary \\  
% Option  
(B) Posterior pituitary \\  
% Option  
(C) Hypothalamus \\  
% Option  
(D) Adrenal gland \\  

\bigskip  

% Correct Answer  
\noindent \textbf{Correct Answer:} (B) Posterior pituitary  

\bigskip  

% Solution  
\noindent \textbf{Solution:} \\  

\textbf{Step 1: Understanding the Question:} \\
The question is based on the human endocrine system, specifically the site of secretion of neurohypophyseal hormones. \\
It requires identifying which gland is responsible for releasing oxytocin and vasopressin into the systemic circulation. \\

\textbf{Step 2: Key Formula or Approach:} \\
This is a concept-based biology question. \\
The approach involves distinguishing between the site of hormone synthesis and the site of hormone secretion. \\

\textbf{Step 3: Detailed Explanation:} \\
\begin{itemize}
\item Oxytocin and vasopressin (or antidiuretic hormone, ADH) are nonapeptides synthesized in the hypothalamus. \\
\item Specifically, oxytocin is synthesized in the neurons of the paraventricular nucleus, while vasopressin is synthesized in the supraoptic nucleus. \\
\item Following synthesis, these hormones are packaged into secretory granules and transported axonally down the hypothalamo-hypophyseal tract. \\
\item They terminate and are stored in the posterior pituitary gland (neurohypophysis) in specialized axon terminals called Herring bodies. \\
\item Upon stimulation by appropriate electrical signals, these hormones are released via exocytosis into the capillary plexus of the posterior pituitary. \\
\item Thus, while the hypothalamus is the site of synthesis, the posterior pituitary is the anatomical site of storage and secretion into the bloodstream. \\
\end{itemize}

\textbf{Step 4: Final Answer:} \\
Since the question asks which organ secretes these hormones, the correct option is the posterior pituitary. \\

\bigskip  

% Quick Tip  
\begin{quicktipbox}  
Remember that the posterior pituitary (neurohypophysis) does not synthesize any hormones of its own. \\
It only stores and secretes hormones that are synthesized by the hypothalamus. \\
\end{quicktipbox}  

% Topic - Endocrine System  
\vspace{0.5cm}  
\hrule  
\vspace{0.5cm}  

\noindent \textbf{2.}  
% Question text  
\textbf{The outer layer of Dicot root is called:} \\  

% Option  
(A) Epidermis \\  
% Option  
(B) Endodermis \\  
% Option  
(C) Pericycle \\  
% Option  
(D) Cortex \\  

\bigskip  

% Correct Answer  
\noindent \textbf{Correct Answer:} (A) Epidermis  

\bigskip  

% Solution  
\noindent \textbf{Solution:} \\  

\textbf{Step 1: Understanding the Question:} \\
The question requires us to identify the outermost tissue layer of a dicotyledonous root. \\
This is a standard question from plant anatomy focusing on root tissue organization. \\

\textbf{Step 2: Key Formula or Approach:} \\
We examine the concentric arrangement of tissue layers in a typical dicot root from the outside to the inside. \\
The order of layers is: Epidermis (Epiblema) $\rightarrow$ Cortex $\rightarrow$ Endodermis $\rightarrow$ Pericycle $\rightarrow$ Vascular Bundles. \\

\textbf{Step 3: Detailed Explanation:} \\
\begin{itemize}
\item The outermost layer of a dicot root is called the epidermis, also commonly referred to as the epiblema or piliferous layer. \\
\item It consists of a single layer of compactly arranged, thin-walled parenchymatous cells. \\
\item Unlike the stem epidermis, the root epidermis typically lacks a thick cuticle and stomata to allow for efficient water absorption. \\
\item Many of the epidermal cells protrude outwards to form unicellular root hairs, which significantly increase the surface area for absorption of water and minerals. \\
\item The layer directly beneath the epidermis is the cortex, which consists of multiple layers of loosely packed parenchyma cells. \\
\item The endodermis is the innermost layer of the cortex, characterized by barrel-shaped cells containing Casparian strips. \\
\item The pericycle lies inward to the endodermis and gives rise to lateral roots. \\
\end{itemize}

\textbf{Step 4: Final Answer:} \\
Thus, the outermost layer of the dicot root is the epidermis (epiblema). \\

\bigskip  

% Quick Tip  
\begin{quicktipbox}  
In root anatomy, the terms 'epidermis', 'epiblema', and 'piliferous layer' are used interchangeably. \\
Look for these synonyms when solving questions about root tissue systems. \\
\end{quicktipbox}  

% Topic - Plant Anatomy  
\vspace{0.5cm}  
\hrule  
\vspace{0.5cm}  

\noindent \textbf{3.}  
% Question text  
\textbf{The notch of kidney where vessels enter is called:} \\  

% Option  
(A) Hilum \\  
% Option  
(B) Pelvis \\  
% Option  
(C) Calyx \\  
% Option  
(D) Ureter \\  

\bigskip  

% Correct Answer  
\noindent \textbf{Correct Answer:} (A) Hilum  

\bigskip  

% Solution  
\noindent \textbf{Solution:} \\  

\textbf{Step 1: Understanding the Question:} \\
The question asks for the anatomical name of the specific notch on the kidney where blood vessels, nerves, and ureters enter or exit. \\
This is a fundamental question of renal anatomy. \\

\textbf{Step 2: Key Formula or Approach:} \\
The approach involves identifying the structural entry and exit portal on the medial surface of the kidney. \\
We analyze the role and position of the hilum, pelvis, calyx, and ureter. \\

\textbf{Step 3: Detailed Explanation:} \\
\begin{itemize}
\item The kidney is a bean-shaped organ located in the retroperitoneal space. \\
\item On the medial concave border of each kidney, there is a deep, vertical slit or notch known as the renal hilum (or hilus). \\
\item The hilum serves as the primary gateway for structures entering and leaving the kidney. \\
\item These structures include the renal artery, renal vein, lymphatic vessels, nerves, and the renal pelvis which leads to the ureter. \\
\item The renal pelvis is a funnel-shaped structure internal to the hilum that collects urine from the calyces. \\
\item Calyces (minor and major) are cup-like subdivisions of the pelvis that receive urine from the renal papillae of the medullary pyramids. \\
\item The ureter is the long muscular tube that carries urine from the renal pelvis to the urinary bladder. \\
\end{itemize}

\textbf{Step 4: Final Answer:} \\
The notch on the medial surface where vessels and nerves enter the kidney is the hilum. \\

\bigskip  

% Quick Tip  
\begin{quicktipbox}  
To easily recall renal landmarks, remember: \\
Hilum = Entrance/Exit Gate. \\
Pelvis = Funnel collector. \\
Ureter = Transport pipe. \\
\end{quicktipbox}  

% Topic - Human Excretory System  
\vspace{0.5cm}  
\hrule  
\vspace{0.5cm}  

\noindent \textbf{4.}  
% Question text  
\textbf{Who is called the Father of Ecology in India?} \\  

% Option  
(A) Ramdeo Misra \\  
% Option  
(B) S. P. Singh \\  
% Option  
(C) Salim Ali \\  
% Option  
(D) M. S. Swaminathan \\  

\bigskip  

% Correct Answer  
\noindent \textbf{Correct Answer:} (A) Ramdeo Misra  

\bigskip  

% Solution  
\noindent \textbf{Solution:} \\  

\textbf{Step 1: Understanding the Question:} \\
The question asks for the identity of the scientist recognized as the "Father of Ecology in India". \\
This is a factual and historical question in the field of Indian environmental biology. \\

\textbf{Step 2: Key Formula or Approach:} \\
We identify the contributions of the listed scientists to find who pioneered ecological research and education in India. \\

\textbf{Step 3: Detailed Explanation:} \\
\begin{itemize}
\item Professor Ramdeo Misra is widely recognized as the Father of Ecology in India. \\
\item He established the first dedicated postgraduate course in ecology at Banaras Hindu University (BHU) in Varanasi. \\
\item His extensive research on tropical communities, primary productivity, and nutrient cycling laid the foundation for modern ecological studies in the country. \\
\item He also played an instrumental role in the establishment of the National Committee for Environmental Planning and Coordination, which later evolved into the Ministry of Environment and Forests. \\
\item Let us evaluate the other options: \\
\item S. P. Singh is a prominent ecologist known for his work on forest ecology, but is not the founder. \\
\item Salim Ali is famously known as the "Birdman of India" and was a renowned ornithologist and naturalist. \\
\item M. S. Swaminathan is internationally celebrated as the "Father of the Green Revolution in India" for his work in agriculture and genetics. \\
\end{itemize}

\textbf{Step 4: Final Answer:} \\
Therefore, Professor Ramdeo Misra is the correct choice as the pioneer of Indian ecology. \\

\bigskip  

% Quick Tip  
\begin{quicktipbox}  
In biological history questions, associate key terms with scientists: \\
Ramdeo Misra = Indian Ecology. \\
M. S. Swaminathan = Green Revolution. \\
Salim Ali = Ornithology. \\
\end{quicktipbox}  

% Topic - Ecology  
\vspace{0.5cm}  
\hrule  
\vspace{0.5cm}  

\noindent \textbf{5.}  
% Question text  
\textbf{Which animal was not Uricotelic?} \\  

% Option  
(A) Birds \\  
% Option  
(B) Land snail \\  
% Option  
(C) Reptiles \\  
% Option  
(D) Adult Frog \\  

\bigskip  

% Correct Answer  
\noindent \textbf{Correct Answer:} (D) Adult Frog  

\bigskip  

% Solution  
\noindent \textbf{Solution:} \\  

\textbf{Step 1: Understanding the Question:} \\
The question asks to identify the animal among the given options that is NOT uricotelic. \\
Uricotelic animals are those that excrete nitrogenous waste primarily in the form of uric acid. \\

\textbf{Step 2: Key Formula or Approach:} \\
We analyze the mode of excretion of different animal groups based on water conservation needs. \\
Animals are classified into three major groups: \\
1. Ammonotelic (excrete ammonia, requires high water). \\
2. Ureotelic (excrete urea, requires moderate water). \\
3. Uricotelic (excrete uric acid, requires minimal water). \\

\textbf{Step 3: Detailed Explanation:} \\
\begin{itemize}
\item Uricotelic organisms excrete uric acid or its salts as a paste or dry pellet. \\
\item Uric acid is highly insoluble in water and relatively non-toxic, allowing organisms to conserve maximum body water. \\
\item This is a vital adaptation for terrestrial environments with limited water availability, such as those inhabited by birds, land snails, insects, and reptiles. \\
\item Consequently, birds, land snails, and reptiles are classically categorized as uricotelic. \\
\item On the other hand, the adult frog is a semi-aquatic amphibian. \\
\item Adult frogs live in environments where water is accessible, and they excrete their nitrogenous waste as urea, which is soluble in water and moderately toxic. \\
\item Thus, adult frogs are ureotelic, not uricotelic. \\
\item Note that the larval stage of a frog (tadpole) is ammonotelic because it lives exclusively in water. \\
\end{itemize}

\textbf{Step 4: Final Answer:} \\
The adult frog is ureotelic, making it the correct option for an animal that is not uricotelic. \\

\bigskip  

% Quick Tip  
\begin{quicktipbox}  
An easy way to remember the excretory products based on habitat: \\
Aquatic (fish, tadpoles) $\rightarrow$ Ammonotelic. \\
Terrestrial/Semi-aquatic (mammals, adult frogs) $\rightarrow$ Ureotelic. \\
Highly dry-adapted/Aerial (birds, reptiles, snails) $\rightarrow$ Uricotelic. \\
\end{quicktipbox}  

% Topic - Excretion  
\vspace{0.5cm}  
\hrule  
\vspace{0.5cm}  
